\chapter{Introducción}
\label{chap:introduccion}

\section{Contextualización}

Rara vez se repara en la red de datos de una universidad mientras funciona; se la advierte, en cambio, en el instante en que una videoconferencia se congela o una plataforma académica tarda en responder. Esta invisibilidad cotidiana encierra una paradoja: cuanto más indispensable se vuelve una infraestructura, menos se la observa y, por tanto, menos se conoce el servicio que efectivamente entrega. Una Institución de Educación Superior sostiene sobre su red de datos la operación de los sistemas académicos, las plataformas de gestión del aprendizaje, los laboratorios de cómputo, la biblioteca digital y el correo institucional. A ella acceden de forma simultánea estudiantes, docentes, investigadores y personal administrativo, en su mayoría desde dispositivos que la institución no gestiona y que se asocian a puntos de acceso inalámbrico. Esta circunstancia sitúa en el extremo cliente la última porción de la ruta y, con ella, una parte sustantiva de la degradación que el usuario percibe.

Ahora bien, sería un error concebir esta red como una red corporativa más. Las redes de campus universitario poseen una fisonomía propia, que la literatura ha descrito con precisión \cite{kotz2005campus,henderson2004changing}. En primer lugar, la demanda no se distribuye de manera uniforme, sino que se concentra en las franjas horarias que coinciden con el inicio de las jornadas de clase. En segundo lugar, la población de dispositivos es heterogénea y, en su mayoría, ajena a la gestión institucional, bajo esquemas de tipo \textit{Bring Your Own Device} (BYOD). Por último, sobre un mismo medio coexisten perfiles de tráfico con exigencias dispares, que van desde la navegación académica hasta la videoconferencia sincrónica y el acceso a sistemas transaccionales. De ello se sigue que el desempeño de una red de campus no es un valor único, sino una magnitud que varía según el lugar y el momento en que se observa.

¿Desde dónde, entonces, debe observarse una red para conocer el servicio que presta? El diagnóstico admite dos planos de observación que responden a preguntas distintas. El plano de gestión de la infraestructura, al que solo se accede con credenciales administrativas, ofrece los contadores de interfaz y la exportación de flujos de los equipos de red; describe, por tanto, la utilización de los enlaces y la composición agregada del tráfico. El plano del extremo, en cambio, no requiere privilegio alguno y cuantifica el servicio efectivamente entregado en el punto de consumo: la capacidad disponible en la ruta, el retardo acumulado por salto y el comportamiento de los protocolos de transporte y de aplicación. Ninguno de los dos sustituye al otro. Un enlace troncal holgado no garantiza que el usuario reciba un buen servicio, del mismo modo que una sesión lenta no prueba que el enlace esté saturado.

A esta distinción técnica se suma una exigencia institucional propia del sistema ecuatoriano de educación superior. El modelo de evaluación externa con fines de acreditación vigente para las universidades y escuelas politécnicas dedica uno de sus indicadores a la infraestructura física y tecnológica, y exige que la institución no solo disponga de ella, sino que la monitoree y evalúe su funcionalidad con miras a la mejora continua \cite{caces2023modelo}. La calidad de la red deja de ser, así, una cuestión exclusivamente operativa para convertirse en un asunto de aseguramiento de la calidad institucional.

De esta doble condición se desprende el propósito del presente trabajo. Si una parte decisiva del desempeño se define en el extremo, y si ese extremo puede observarse sin intervenir sobre la infraestructura, la evidencia sobre el servicio entregado puede obtenerse desde el propio equipo de un usuario autorizado. Es en este contexto donde la carrera de Ingeniería en Tecnologías de la Información plantea la caracterización instrumentada del desempeño extremo a extremo de la red de una Institución de Educación Superior desde el segmento de acceso del cliente. Para ello recurre a dos técnicas de naturaleza distinta y complementaria: la medición activa, que inyecta flujos de prueba controlados para derivar los parámetros de capacidad y de retardo de la ruta, y la disección pasiva de las tramas del propio equipo, que permite atribuir ese retardo a cada fase del establecimiento del servicio.

\section{Estado del Arte}

¿Qué sabe ya la disciplina sobre el desempeño que una red entrega a sus usuarios, y desde dónde lo ha observado? Responder a esta pregunta exige recorrer una tradición de investigación que ha medido las redes desde lugares muy distintos: desde los equipos de la infraestructura, desde los puntos de acceso, desde la pasarela del hogar o desde el propio equipo del usuario. El recorrido que sigue no pretende ser exhaustivo; busca, más bien, mostrar cómo cada punto de observación ilumina una faceta del problema y deja, al mismo tiempo, otra en penumbra.

El punto de partida obligado lo constituye el trabajo de Kotz y Essien \cite{kotz2005campus}, uno de los primeros estudios a gran escala sobre una red inalámbrica de campus en producción. En Dartmouth College, sobre una red de 476 puntos de acceso distribuidos en 161 edificaciones, los autores siguieron durante once semanas la actividad de cerca de dos mil usuarios y combinaron registros del sistema, consultas SNMP y capturas de tráfico para caracterizar los patrones de uso, la movilidad y la distribución del tráfico por aplicación. Su observación se sitúa, por tanto, en el plano de gestión: describe con riqueza cómo se usa la red, pero no qué servicio recibe cada usuario en términos de capacidad, retardo o pérdida.

Dos décadas después, ese mismo plano de observación conserva su vigencia en el contexto latinoamericano. Al aplicar analítica de datos sobre los registros de la controladora inalámbrica de la Universidad Nacional del Altiplano, en Perú, Cari Quispe \cite{cari2025altiplano} identificó picos de uso entre las 10:00 y las 14:00 y puntos de acceso sobrecargados, definidos como aquellos cuya latencia media supera los 40~ms y cuya pérdida de paquetes excede el 1,5\,\%. Su estudio confirma la concentración horaria de la demanda propia de las redes universitarias y ofrece un antecedente de localización de la degradación en el espacio y en el tiempo. Sus umbrales, sin embargo, son propios y no se derivan de una recomendación técnica, y su fuente de datos exige, una vez más, acceso a la infraestructura.

Si los trabajos anteriores describen el uso, otros se proponen medir directamente el retardo. Sui \textit{et al.} \cite{sui2016wifiseer} desplegaron en el campus de la Universidad Tsinghua un sistema que asocia mediciones de latencia del tramo inalámbrico con once factores de radio, tres de protocolo y seis de entorno, obtenidos por SNMP desde la infraestructura, sobre más de 2700 puntos de acceso, 114 edificaciones y más de 47\,000 dispositivos. Encontraron que esa latencia sigue una distribución de cola larga ---con una mediana de apenas 3~ms, pero percentiles 90 y 99 cercanos a 20~ms y 250~ms--- y que la utilización del canal la explica mejor que la intensidad de la señal recibida. Dos lecciones se desprenden para el presente trabajo. La primera es que los promedios ocultan precisamente la degradación que el usuario percibe, de modo que una línea base debe describir la dispersión y no solo la tendencia central. La segunda es metodológica: los autores advierten que el \textit{ping} convencional emitido desde un dispositivo móvil se ve inflado por los modos de ahorro de energía, efecto que cualquier medición desde el extremo debe controlar.

En el mismo campus, Pei \textit{et al.} \cite{pei2016weakest} descompusieron el retardo de ida y vuelta de los paquetes TCP en sus tramos inalámbrico de subida, cableado e inalámbrico de bajada, a partir de 47 puntos de acceso observados durante dos meses. Comprobaron que el tramo inalámbrico puede ser el eslabón más débil de la ruta: más del 50\,\% de los paquetes TCP sufrió en él un retardo superior a 20~ms, el 10\,\% uno superior a 100~ms, y en más de la mitad de los casos ese tramo concentró más del 60\,\% del retardo de ida y vuelta. Con un árbol de decisión que relaciona ese retardo con factores como la intensidad de la señal del cliente y la tasa de reintentos, lograron en un caso real reducir la mediana del retardo de 50~ms a 10~ms. El estudio prueba que la contribución del acceso inalámbrico puede aislarse y atribuirse. Su método no exige instrumentar al cliente, pero sí instalar la herramienta en puntos de acceso con \textit{firmware} OpenWrt; es decir, intervenir sobre la infraestructura.

Medir el retardo, con todo, no basta para explicar por qué una transferencia es lenta. Zhang \textit{et al.} \cite{zhang2002flowrates} abordaron esta cuestión a partir de trazas de paquetes capturadas en varios enlaces de internet y desarrollaron una herramienta, T-RAT, que analiza la dinámica de cada conexión TCP para inferir el factor que limita su tasa de transmisión. En sus trazas, las causas más frecuentes resultaron ser la congestión de la red y el límite impuesto por la ventana del receptor. El trabajo demuestra que la disección pasiva de las sesiones de transporte permite distinguir limitaciones de naturaleza muy distinta, que exigen a su vez remedios distintos; su punto de observación son, sin embargo, los enlaces de la red y no el equipo del usuario.

La misma lógica de atribución puede extenderse a las fases que preceden a la transferencia. Sundaresan \textit{et al.} \cite{sundaresan2013web} midieron el tiempo de carga de nueve sitios web populares desde más de 5000 redes de acceso de banda ancha y comprobaron que, cuando el caudal descendente del acceso supera unos 16~Mbit/s, la latencia pasa a ser el principal cuello de botella del tiempo de carga. Mediante una herramienta instalada en la pasarela del hogar, Mirage, descompusieron ese tiempo en la resolución DNS, el establecimiento de la conexión TCP y la descarga de los objetos, y mostraron que optimizaciones sencillas de almacenamiento en caché pueden reducirlo hasta en un 53\,\%. La conclusión es directamente trasladable al presente trabajo: más capacidad no siempre significa mejor servicio, y solo la descomposición por fases permite saber cuándo el problema reside en otra parte.

Existe, además, una fuente de degradación que rara vez se examina porque no deja rastro visible. Luckie y Stasiewicz \cite{luckie2010pmtud} midieron el comportamiento del mecanismo de descubrimiento de la unidad máxima de transmisión de la ruta en 50\,000 sitios web populares y encontraron tasas de falla de entre el 5\,\% y el 18\,\%, según la unidad máxima de transmisión del enlace restrictivo; observaron también que el 11\,\% de los operadores limita sus paquetes a 1380 bytes para no depender de ese mecanismo. Cuando este falla, las comunicaciones se interrumpen de un modo difícil de diagnosticar, lo que justifica incluir su verificación explícita en una caracterización desde el extremo.

El contexto ecuatoriano ofrece, por su parte, el antecedente más próximo en cuanto a instrumentación. Moreno-Arteaga \textit{et al.} \cite{moreno2023sitesurvey} estudiaron la red inalámbrica de acceso libre del campus Lodana de la Universidad Técnica de Manabí, motivados expresamente por el peso de la conectividad en la evaluación institucional del CACES. Realizaron estudios de sitio con la herramienta Acrylic Wi-Fi HeatMaps desde un equipo cliente, apoyados en un servidor iperf3 instalado temporalmente en la red de área local de dos facultades, y contrastaron los resultados con umbrales de $-75$~dBm de intensidad de señal, 20~dB de relación señal-ruido, 200~ms de tiempo de ida y vuelta y 2\,\% de pérdida de paquetes; propusieron luego dos diseños alternativos, evaluados por simulación en OPNET. El trabajo valida, en una universidad del país, la medición desde un equipo cliente hacia un destino interno a la red. Sus umbrales, sin embargo, no proceden de una recomendación técnica, y su alcance se limita a la banda de 2,4~GHz.

Queda por establecer en qué medida la concentración de la demanda degrada efectivamente el desempeño medido. Toledo González y Ríos Gómez \cite{toledo2026densidad} lo cuantificaron en un laboratorio controlado, con Wireshark e iperf3, bajo tres escenarios de densidad de dispositivos en una red de área local cuyo medio de transmisión no se especifica: al pasar de baja a alta densidad, la latencia aumentó de 1,30~ms a 4,80~ms, la pérdida de paquetes de 0,5\,\% a 5,8\,\%, y el ancho de banda disponible descendió de 250~Mbps a 90~Mbps. Aunque el entorno experimental no reproduce la complejidad de una red en producción, estos resultados justifican que el diseño muestral del presente trabajo cubra tanto las franjas de alta como las de baja ocupación.

La Tabla~\ref{tab:estado_arte} reúne lo que el recorrido anterior ha puesto de manifiesto. En ella, cada trabajo se contrasta con los criterios que, en correspondencia con los objetivos específicos de este proyecto, definen una caracterización del desempeño desde el extremo: el punto desde el que se observa la red, la medición activa, la disección pasiva de protocolos, la descomposición del retardo por fase, el registro de los indicadores de capa física y de enlace, el muestreo por lugar y franja horaria y el contraste con umbrales normativos.

\begin{table}[htbp]
  \centering
  \caption{Cobertura de los trabajos analizados frente a los criterios que definen el alcance del proyecto.}
  \label{tab:estado_arte}
  \footnotesize
  \renewcommand{\arraystretch}{1.05}
  \setlength{\tabcolsep}{4pt}
  \newcommand{\girar}[1]{\rotatebox{90}{\textbf{#1}\hspace{6pt}}}
  \newcommand{\contexto}[1]{\newline{\scriptsize\itshape #1}}
  \begin{tabularx}{\textwidth}{|>{\raggedright\arraybackslash}p{3.1cm}|>{\raggedright\arraybackslash}X|c|c|c|c|c|c|}
    \hline
    & & \multicolumn{2}{c|}{\scriptsize\textbf{Técnica}} & \multicolumn{4}{c|}{\scriptsize\textbf{Análisis}} \\
    \cline{3-8}
    \textbf{Trabajo y contexto} & \textbf{Punto de observación}
      & \girar{Medición activa} & \girar{Disección pasiva}
      & \girar{Desglose por fase} & \girar{Capa física y enlace} & \girar{Lugar y franja} & \girar{Umbrales normativos} \\
    \hline
    Zhang \textit{et al.} \cite{zhang2002flowrates}\contexto{Enlaces de internet, EE.\,UU.}
      & Trazas de paquetes en enlaces de la red
      & \NoCumple & \Cumple & \NoCumple & \NoCumple & \NoCumple & \NoCumple \\
    \hline
    Kotz y Essien \cite{kotz2005campus}\contexto{Dartmouth College, EE.\,UU.}
      & Plano de gestión: \textit{syslog}, SNMP y captura de tráfico
      & \NoCumple & \Parcial & \NoCumple & \NoCumple & \Parcial & \NoCumple \\
    \hline
    Luckie y Stasiewicz \cite{luckie2010pmtud}\contexto{50\,000 sitios web}
      & Equipo de medición propio con sondas activas
      & \Cumple & \Parcial & \NoCumple & \NoCumple & \NoCumple & \NoCumple \\
    \hline
    Sundaresan \textit{et al.} \cite{sundaresan2013web}\contexto{Más de 5000 redes de acceso}
      & Pasarela del hogar
      & \Cumple & \NoCumple & \Cumple & \NoCumple & \NoCumple & \NoCumple \\
    \hline
    Sui \textit{et al.} \cite{sui2016wifiseer}\contexto{U. Tsinghua, China}
      & Infraestructura: sondeos de latencia y SNMP
      & \Cumple & \NoCumple & \NoCumple & \Cumple & \Parcial & \NoCumple \\
    \hline
    Pei \textit{et al.} \cite{pei2016weakest}\contexto{U. Tsinghua, China}
      & Puntos de acceso con \textit{firmware} OpenWrt
      & \NoCumple & \Cumple & \Parcial & \Cumple & \Parcial & \Parcial \\
    \hline
    Cari Quispe \cite{cari2025altiplano}\contexto{U. Nacional del Altiplano, Perú}
      & Plano de gestión: registros de la controladora inalámbrica
      & \NoCumple & \NoCumple & \NoCumple & \NoCumple & \Cumple & \Parcial \\
    \hline
    Moreno-Arteaga \textit{et al.} \cite{moreno2023sitesurvey}\contexto{U. Técnica de Manabí, Ecuador}
      & Equipo cliente con servidor iperf3 interno
      & \Cumple & \NoCumple & \NoCumple & \Cumple & \Parcial & \Parcial \\
    \hline
    Toledo y Ríos \cite{toledo2026densidad}\contexto{Laboratorio controlado}
      & Equipos de prueba en laboratorio
      & \Cumple & \Parcial & \NoCumple & \NoCumple & \NoCumple & \NoCumple \\
    \hline\hline
    \textbf{Presente proyecto}\contexto{Institución de Educación Superior, Ecuador}
      & Equipo del usuario, sin privilegios administrativos
      & \Cumple & \Cumple & \Cumple & \Cumple & \Cumple & \Cumple \\
    \hline
  \end{tabularx}
  \vspace{2pt}
  \raggedright
  \elaboradopor{Josué Gabriel Godoy Apolo}
  \fuente{Elaboración propia con base en los trabajos citados}
  \noindent\textit{Nota:} \Cumple{} aborda el criterio; \Parcial{} lo aborda de forma parcial (captura tráfico sin disecar todas las sesiones, descompone el retardo por tramos de la ruta y no por fases del servicio, distingue solo tipos de lugar o de escenario, o aplica umbrales propios sin sustento normativo); \NoCumple{} no lo aborda. \textit{Capa física y enlace} designa el registro, junto con el desempeño, de indicadores como la intensidad de señal, la relación señal-ruido, la modulación o los reintentos; \textit{umbrales normativos}, el contraste de los resultados con los objetivos de una recomendación técnica.
\end{table}

La lectura de la tabla arroja un balance inequívoco: el conocimiento disponible está repartido entre puntos de observación que no se comunican. Los trabajos que relacionan el desempeño con la capa física y de enlace ---Sui \textit{et al.}, Pei \textit{et al.} y Moreno-Arteaga \textit{et al.}--- lo hacen desde la infraestructura o con instrumental de estudio de sitio, mientras que los que descomponen o atribuyen la degradación ---Zhang \textit{et al.} y Sundaresan \textit{et al.}--- observan la red desde sus enlaces o desde la pasarela del hogar, lejos de un campus universitario. Ninguno combina la medición activa con la disección pasiva desde el equipo del usuario. Y ninguno, por último, contrasta sus resultados con los objetivos de una recomendación técnica: los umbrales que aparecen son, en todos los casos, propios.

Entre estos vacíos ---el de una observación desde el extremo que no se limite a medir y el de un contraste con referentes normativos--- se abre un espacio que ninguno de los trabajos ocupa por completo. Es precisamente ahí donde se sitúa el presente proyecto: caracterizar el desempeño extremo a extremo desde el equipo de un usuario autorizado, sin acceso a la infraestructura, mediante la combinación de la medición activa y la disección pasiva, con la descomposición del retardo por fase, la correlación con los indicadores del acceso inalámbrico, un muestreo por punto de asociación y franja horaria y un contraste explícito con los objetivos de la Recomendación ITU-T Y.1541.

\section{Planteamiento del Problema}

Cuando un usuario afirma que «la red está lenta», enuncia un síntoma, no un diagnóstico. La lentitud percibida puede originarse en una capacidad insuficiente de la ruta, pero también en la resolución de nombres, en la negociación de una sesión cifrada, en el control de congestión del transporte o en la calidad del enlace inalámbrico que asocia el dispositivo a la red. Sin evidencia medida desde el extremo, estas causas se confunden y la degradación tiende a imputarse, por defecto, a la capacidad de la red. El problema que aborda este trabajo puede formularse, entonces, como una pregunta: ¿qué desempeño entrega efectivamente la red de datos de una Institución de Educación Superior en el extremo cliente, y a qué componentes de la ruta y de los protocolos puede atribuirse su degradación?

Responder a esta pregunta exige atender varios aspectos técnicos, que se articulan en cuatro planos. El primero es el de la capacidad y el retardo: es preciso determinar el \textit{goodput} sostenido disponible en la ruta, conforme a una metodología de medición normalizada, y cuantificar el retardo de ida y vuelta, la variación del retardo de paquete y la tasa de pérdida, parámetros que condicionan el desempeño de los servicios sincrónicos de voz y video. Tales valores, además, no son constantes: su dispersión según el punto de asociación y la franja horaria forma parte del problema. El segundo plano es el de la atribución. Para atribuir correctamente la degradación observada, el retardo percibido debe descomponerse en sus fases de resolución de nombres, establecimiento de sesión TCP, negociación TLS y transferencia, y las sesiones de transporte deben examinarse en busca de retransmisiones, reordenamiento y limitación por ventana de congestión.

El tercer plano corresponde al acceso inalámbrico, cuya contribución al desempeño observado debe caracterizarse a partir de los indicadores de señal, modulación y reintentos que expone el propio dispositivo. El cuarto, por último, atañe a la ruta misma: la unidad máxima de transmisión efectiva y la operatividad del mecanismo que la descubre, cuyo mal funcionamiento puede detener sesiones de forma silenciosa. A estos cuatro planos se añade una exigencia de método. El procedimiento de medición debe ser instrumentado, reproducible y ejecutable sin privilegios administrativos sobre la infraestructura, de modo que pueda repetirse en el tiempo y compararse consigo mismo.

Las consecuencias de dejar sin respuesta estas cuestiones se manifiestan en varios planos. En el técnico, una degradación mal atribuida conduce a remedios que no corresponden a la causa: ampliar el ancho de banda contratado no corrige una resolución de nombres lenta ni una tasa elevada de reintentos en el enlace inalámbrico. En el de la gestión, sin una localización de la degradación en el espacio y en el tiempo, la institución carece de criterio para priorizar su intervención. Y, en el tiempo, sin una línea base reproducible se vuelve imposible verificar si el servicio que recibe la comunidad universitaria mejora o se deteriora.

\section{Justificación}

Cabe preguntarse, entonces, por qué medir desde el extremo, cuando la infraestructura podría observarse desde su plano de gestión. La respuesta que aquí se sostiene es que la caracterización desde el extremo no es un sustituto empobrecido de la observación administrativa, sino una línea de trabajo autónoma, que responde a una pregunta que el plano de gestión no puede contestar. Esta tesis puede defenderse desde tres perspectivas complementarias.

Desde la perspectiva técnica, el \textit{goodput} y el retardo observados en el extremo no se deducen de la utilización de los enlaces. Dependen del comportamiento del control de congestión del transporte, de la relación entre el tamaño de la ventana y el tiempo de ida y vuelta, de la pérdida acumulada en la ruta \cite{rfc6349}, de la unidad máxima de transmisión efectiva y de la tasa de reintentos del acceso inalámbrico. La disección de las sesiones permite, además, distinguir una degradación originada en la resolución de nombres o en la negociación TLS de una limitación real de capacidad, atribución que condiciona cualquier decisión correctiva posterior. Toda optimización parte, en efecto, de una caracterización correcta del estado de la red \cite{oppenheimer2011topdown}.

Pero el valor de esta caracterización no se agota en el plano técnico. Desde la perspectiva de la gestión, el trabajo aporta evidencia medida sobre el desempeño percibido por la comunidad universitaria, obtenida sin intervención sobre la infraestructura en operación y, por tanto, sin riesgo para la continuidad del servicio. El diseño muestral por punto de asociación y franja horaria permite localizar la degradación en el espacio y en el tiempo, y priorizar en consecuencia. Esta evidencia responde, además, a una exigencia expresa del sistema de aseguramiento de la calidad: el modelo de evaluación del CACES requiere que las instituciones monitoreen y evalúen la funcionalidad de su infraestructura tecnológica \cite{caces2023modelo}. El protocolo de medición documentado permitirá repetir la caracterización de forma periódica, de modo que el estudio no se reduzca a una fotografía aislada, sino que inaugure una práctica sostenida y verificable.

A estas razones se añade, por último, una justificación académica. El trabajo concentra competencias centrales del área de redes de la carrera en un alcance ejecutable por un estudiante en diez semanas: el diseño de un protocolo de medición normalizado, la instrumentación en el extremo, la disección de protocolos por capas, la interpretación de parámetros de calidad de servicio frente a las recomendaciones de la Unión Internacional de Telecomunicaciones y la elaboración de un informe sustentado en evidencia. Y lo hace sobre una red en producción, lo que confiere a sus resultados una validez que ningún escenario de laboratorio podría ofrecer.

¿Quiénes se benefician de este esfuerzo? De manera directa, la institución y el personal técnico responsable de la operación y la planificación de su infraestructura de red, que recibe evidencia técnica instrumentada obtenida sin intervención sobre los equipos en operación. De manera indirecta, la comunidad universitaria en su conjunto ---estudiantes, docentes, investigadores y personal administrativo---, que utiliza los servicios soportados sobre la red.

Resta, con todo, una objeción posible: la de la viabilidad. El proyecto es factible en el plano técnico, ya que recurre exclusivamente a herramientas gratuitas, de código abierto o incluidas en el sistema operativo, como iperf3, mtr, fping, tracepath, dig, curl, Wireshark y tcpdump, y a un equipo portátil con interfaces cableada e inalámbrica. Lo es también en el plano operativo, porque no depende de autorizaciones de acceso a la infraestructura: todas las actividades se ejecutan desde el equipo del estudiante, con las mismas credenciales y privilegios de cualquier miembro de la comunidad universitaria. Esta condición no es una limitación sobrevenida, sino un criterio de diseño del alcance. Lo es, por último, en el plano ético, puesto que la captura se restringe a las interfaces del propio equipo y los flujos de prueba se dimensionan para no degradar el medio compartido.

En suma, el proyecto trasciende el interés de la propia institución. Al aportar evidencia instrumentada sobre el desempeño de la infraestructura tecnológica de una institución de educación superior, se vincula con el Objetivo de Desarrollo Sostenible 9 (Industria, innovación e infraestructura); al caracterizar las condiciones de conectividad en que se desarrolla la actividad académica, se inscribe también en el Objetivo 4 (Educación de calidad).

\section{Objetivos}

\subsection{Objetivo General}

Caracterizar el desempeño extremo a extremo de la red de datos de una Institución de Educación Superior desde el segmento de acceso del cliente, mediante medición activa con iperf3 y mtr, conforme a la Recomendación ITU-T Y.1540 y a la metodología de \textit{goodput} de la RFC 6349, y disección pasiva de protocolos TCP/IP con Wireshark y tcpdump, para la construcción de una línea base de desempeño y de una matriz de hallazgos contrastada con los objetivos de la Recomendación ITU-T Y.1541, sin acceso al plano de gestión de la infraestructura institucional.

\subsection{Objetivos Específicos}

Este propósito general se descompone en cinco objetivos específicos, que reproducen el orden de la caracterización: diseñar el procedimiento, medir la ruta, diseccionar lo capturado, correlacionarlo con el acceso inalámbrico y, solo entonces, evaluar lo obtenido.

\begin{enumerate}
  \item Diseñar el protocolo de medición instrumentada desde el extremo cliente, con los puntos de asociación y las franjas horarias de muestreo, los parámetros de desempeño a registrar, las herramientas y los destinos de referencia, el número de repeticiones por sesión y los criterios de descarte de mediciones anómalas.
  \item Medir, conforme al diseño muestral, el \textit{goodput} sostenido según la metodología de la RFC 6349, el retardo de ida y vuelta, la variación del retardo de paquete, la tasa de pérdida y la unidad máxima de transmisión efectiva de la ruta.
  \item Diseccionar las tramas capturadas en las interfaces del propio equipo, con la descomposición del retardo percibido en sus fases de resolución DNS, establecimiento de sesión TCP, negociación TLS y transferencia, y con la identificación de retransmisiones, reordenamiento y limitación por ventana de congestión.
  \item Correlacionar los parámetros de desempeño obtenidos con los indicadores de capa física y de enlace del acceso inalámbrico registrados en el cliente: potencia de señal recibida, relación señal-ruido, esquema de modulación y codificación negociado y tasa de reintentos.
  \item Evaluar los resultados frente a los objetivos de desempeño de la Recomendación ITU-T Y.1541 y las tasas nominales de las enmiendas IEEE 802.11 en operación, mediante una matriz de hallazgos clasificada por criticidad, base de las recomendaciones técnicas derivadas.
\end{enumerate}

\section{Alcance del Proyecto}

Toda caracterización se define tanto por lo que observa como por el lugar desde el cual observa. El presente trabajo abarca el ciclo completo de la caracterización desde el extremo, desde el diseño del protocolo de medición hasta la entrega del informe, y todas sus actividades se ejecutan desde el equipo del estudiante en su condición de usuario autorizado de la red. Los puntos de asociación y las franjas horarias que serán objeto de estudio se delimitarán en la fase de preparación, en función de los espacios de uso común accesibles al estudiante.

\subsection*{Inclusiones}

El estudio comienza por fijar el modo en que se medirá. Comprende, en consecuencia, el diseño y la documentación del protocolo de medición: los puntos de asociación, la distribución de las franjas horarias, los parámetros de desempeño, las herramientas, los destinos de referencia ---públicos e institucionales, para acotar la contribución de la ruta externa---, el número de repeticiones y los criterios de validez de cada sesión.

Definido el protocolo, el alcance incluye la medición activa de la ruta. El \textit{goodput} sostenido se determina mediante flujos de prueba TCP conforme a la metodología de la RFC 6349, con registro de la ventana de recepción y de la fase de arranque descartada. El retardo de ida y vuelta, la variación del retardo de paquete y la tasa de pérdida se obtienen mediante flujos de prueba UDP y sondas ICMP sobre la ruta de salida. A ello se suman el descubrimiento de la unidad máxima de transmisión efectiva de la ruta, con la verificación de la operatividad del mecanismo de descubrimiento asociado, y el registro de los tiempos de resolución DNS y de los tiempos por fase de carga de servicios institucionales y externos representativos.

Cada medición se acompaña del registro de los indicadores de capa física y de enlace del acceso inalámbrico que expone el cliente: la potencia de señal recibida, la relación señal-ruido, el canal y el ancho de canal asociados, el esquema de modulación y codificación negociado y la tasa de reintentos. En paralelo, se capturan y diseccionan las tramas de las interfaces del propio equipo, a fin de caracterizar el intercambio DHCP, la resolución DNS, el establecimiento de sesión TCP, la negociación TLS y las anomalías de transporte observadas.

El alcance culmina, finalmente, en el análisis. Comprende la construcción de la línea base de desempeño por punto de asociación y franja horaria, con sus medidas de tendencia central y de dispersión; la correlación entre los parámetros de desempeño de capas superiores y los indicadores de capa física y de enlace; el contraste de los resultados con los objetivos de la Recomendación ITU-T Y.1541 y con las tasas nominales de las enmiendas IEEE 802.11 en operación; y la elaboración de la matriz de hallazgos, clasificada por criticidad e impacto sobre la experiencia del usuario, con las recomendaciones técnicas derivadas. Incluye, asimismo, la entrega del protocolo de medición documentado para su ejecución periódica y la socialización de los resultados con el docente tutor y la carrera.

\subsection*{Exclusiones}

Tan importante como lo anterior es precisar lo que el trabajo no hace, y la primera frontera es la del plano de observación. Quedan excluidos el acceso al plano de gestión de los equipos de la infraestructura institucional, la lectura de sus configuraciones y la consulta de contadores mediante protocolos de administración. Por la misma razón, el trabajo no contempla la configuración de puertos espejo, la derivación pasiva de tráfico ni la recolección de flujos exportados por los equipos de red, como tampoco la caracterización de la composición agregada del tráfico institucional, que exigiría precisamente esa observación. El trabajo no requiere, por tanto, autorización del área responsable de la infraestructura.

Una segunda frontera es de orden ético. La captura, la inspección y la disección se restringen al tráfico de las interfaces del propio equipo del estudiante, con exclusión de todo tráfico de terceros. Quedan fuera, igualmente, el descubrimiento, la enumeración o el escaneo de \textit{hosts}, puertos o servicios de la red institucional, así como las pruebas de intrusión, el análisis de vulnerabilidades, el descifrado de tráfico protegido y cualquier actividad de explotación sobre los sistemas institucionales.

Por último, otras exclusiones responden a límites de competencia y de medios. Dado su carácter descriptivo, el trabajo no interviene sobre la configuración de la red ni aplica las recomendaciones que formula. No comprende la certificación de enlaces de cableado estructurado ni el diagnóstico de los cuartos de telecomunicaciones o del equipamiento activo, que requieren un acceso físico y administrativo no disponible, ni estudios de sitio inalámbrico con analizadores de espectro profesionales: el registro se limita a los indicadores que expone el controlador de la interfaz del cliente. Tampoco contempla la adquisición de equipamiento, instrumentación especializada o licencias de software comercial, ni la evaluación contractual de los servicios contratados al proveedor de acceso a internet.

Estas fronteras dibujan, en conjunto, el objeto preciso del estudio: la caracterización describe el desempeño entregado en el extremo y no el estado interno de la infraestructura, delimitación que se declarará de forma explícita al interpretar los resultados.

\section{Fundamentación Teórica}

Ninguna medición es neutral: toda medición presupone un modelo de lo que se mide y un criterio de lo que se considera aceptable. Los apartados que siguen exponen ese doble fundamento. Se presentan, en primer lugar, el lugar desde el que se observa la red y las técnicas con que se la observa; luego, los parámetros que describen su desempeño y los protocolos que lo determinan; y, finalmente, los marcos que permiten organizar el trabajo y jerarquizar lo que se encuentre.

\subsection{Planos de observación de una red}

Toda medición de red se realiza desde algún lugar, y ese lugar determina lo que puede conocerse. En el \textbf{plano de gestión}, la información procede de los propios equipos de la infraestructura. El protocolo simple de administración de red (SNMP) permite consultar sus contadores a partir de bases de información de gestión, como la MIB-II definida en la RFC 1213, que reúne los objetos de interfaces, tráfico y errores necesarios para calcular la utilización de los enlaces \cite{rfc1213}; el protocolo IPFIX, definido en la RFC 7011, permite por su parte exportar información de flujos para caracterizar la composición del tráfico por protocolo y por aplicación \cite{rfc7011}. Ambos exigen credenciales administrativas. En el \textbf{plano del extremo}, en cambio, la información procede del equipo de un usuario que intercambia tráfico con destinos de referencia, internos o externos a la institución. Desde allí no puede conocerse la utilización de un enlace ---la propia Recomendación ITU-T Y.1540 advierte que la capacidad utilizada, la utilización y la capacidad disponible no son medibles externamente y requieren sistemas de monitoreo con acceso a los enlaces \cite{ituy1540}---, pero sí el servicio que la ruta completa entrega en el punto de consumo. La Figura~\ref{fig:planos_observacion} representa ambos planos sobre una misma ruta.

\begin{figure}[htbp]
  \centering
  % Planos de observación de una red: gestión de la infraestructura frente a extremo cliente
\begin{tikzpicture}[x=1cm, y=1cm,
  nodo/.style={caja, minimum width=1.7cm, minimum height=0.95cm},
  gestion/.style={-{Stealth[length=1.6mm]}, densely dashed, gray!70!black},
  medida/.style={{Stealth[length=2mm]}-{Stealth[length=2mm]}, thick, FigAcento},
  rot/.style={font=\scriptsize}]
  % Ruta extremo a extremo
  \node[nodo, fill=white, very thick] (cli) at (0,0) {Equipo\\cliente};
  \node[nodo] (ap) at (2.35,0) {Punto de\\acceso};
  \node[nodo] (campus) at (4.7,0) {Red de\\campus};
  \node[nodo] (borde) at (7.05,0) {Borde y\\proveedor};
  \node[nodo, ellipse, minimum width=1.9cm, inner sep=1pt] (inet) at (9.4,0) {Internet};
  \node[nodo, fill=white] (srv) at (11.75,0) {Servidor de\\referencia};
  \draw[FigAcento, thick, densely dotted] (cli) -- (ap);
  \node[rot] at (1.175,0.65) {802.11};
  \draw[FigAcento, thick] (ap) -- (campus) -- (borde) -- (inet) -- (srv);
  % Plano de gestión: los equipos informan a una estación con credenciales
  \node[caja, minimum width=3.2cm, minimum height=0.8cm] (gest) at (4.7,2.1) {Estación de gestión};
  \draw[gestion] (ap.north) -- ($(gest.south)+(-1.1,0)$);
  \draw[gestion] (campus.north) -- (gest.south);
  \draw[gestion] (borde.north) -- ($(gest.south)+(1.1,0)$);
  \node[rot, text=gray!70!black] at (7.4,1.35) {SNMP, IPFIX};
  \node[rot, anchor=south west, text width=13.2cm, align=left] at (-0.85,2.6)
    {\textbf{\textcolor{FigAcento}{Plano de gestión}} (requiere credenciales administrativas): contadores de interfaz y exportación de flujos de los equipos; describe la utilización de los enlaces y la composición agregada del tráfico.};
  % Plano del extremo: el equipo del usuario mide hacia destinos internos y externos
  \draw[medida] ($(cli.south)+(0.3,0)$) |- (4.7,-1.0) -- (campus.south);
  \node[rot, above] at (2.5,-1.0) {Destinos institucionales};
  \draw[medida] ($(cli.south)+(-0.3,0)$) |- (11.75,-1.6) -- (srv.south);
  \node[rot, above] at (6.0,-1.6) {Destinos externos};
  \node[rot, anchor=north west, text width=13.2cm, align=left] at (-0.85,-1.85)
    {\textbf{\textcolor{FigAcento}{Plano del extremo}} (sin privilegios): medición activa y disección pasiva desde el equipo del usuario; cuantifica el servicio entregado en el punto de consumo.};
\end{tikzpicture}

  \caption{Planos de observación de una red: plano de gestión y plano del extremo.}
  \label{fig:planos_observacion}
  \raggedright
  \elaboradopor{Josué Gabriel Godoy Apolo}
  \fuente{Elaboración propia con base en las RFC 1213 y 7011 y en la Recomendación ITU-T Y.1540 \cite{rfc1213,rfc7011,ituy1540}}
\end{figure}

\subsection{Medición activa y disección pasiva}

Desde el extremo, la red puede interrogarse de dos maneras. La \textbf{medición activa} inyecta tráfico de prueba controlado y deriva los parámetros de desempeño de su comportamiento: flujos TCP para estimar la capacidad que la ruta ofrece a una transferencia, flujos UDP para medir la pérdida y la variación del retardo, y sondas ICMP para medir el tiempo de ida y vuelta, salto por salto, hasta el destino. Su virtud es el control, pues se conoce exactamente qué se envió, cuándo y hacia dónde; su costo es que añade carga a un medio compartido, por lo que su volumen debe dimensionarse y registrarse. La \textbf{medición pasiva}, por su parte, no genera tráfico: captura las tramas que atraviesan una interfaz y las disecciona por capas. Cuando la captura se restringe a las interfaces del propio equipo, permite reconstruir, para cada sesión del usuario, el intercambio DHCP, la resolución de nombres, el establecimiento de la conexión TCP, la negociación TLS y las anomalías de transporte, sin observar el tráfico de terceros. Las dos técnicas no compiten entre sí: la primera dice cuánto ofrece la ruta; la segunda, en qué se emplea el tiempo de cada sesión.

\subsection{Parámetros de transferencia de paquetes IP}

Medir el desempeño supone, ante todo, ponerse de acuerdo sobre qué se entiende por desempeño. La Recomendación ITU-T Y.1540 define los parámetros de la transferencia de paquetes IP entre dos puntos de medición \cite{ituy1540}. El \textbf{retardo de transferencia de paquetes IP} (IPTD) es el tiempo que transcurre entre el paso de un paquete por el punto de medición de origen y su paso por el de destino; se trata, por definición, de un retardo en un solo sentido, que puede resumirse por su media, su mínimo o su mediana. La \textbf{variación del retardo de paquetes IP} (IPDV) de dos puntos expresa, para cada paquete, la diferencia entre su retardo y un retardo de referencia, habitualmente el mínimo de la población observada. La \textbf{tasa de pérdida de paquetes IP} (IPLR) es la proporción de paquetes transmitidos que no llegan a su destino, y la \textbf{tasa de paquetes IP reordenados} (IPRR), la proporción de los que llegan fuera de secuencia. El IETF, a través de su grupo de trabajo de métricas de desempeño IP, ha definido métricas equivalentes para el retardo en un sentido \cite{rfc7679}, el retardo de ida y vuelta \cite{rfc2681}, la variación del retardo \cite{rfc3393} y la pérdida \cite{rfc7680}.

Ahora bien, estas definiciones no siempre coinciden con lo que reportan las herramientas. Desde un único extremo, sin relojes sincronizados en el origen y en el destino, no puede medirse el retardo en un solo sentido, sino el tiempo de ida y vuelta (RTT), que lo acota por exceso. Del mismo modo, la variación del retardo admite al menos dos lecturas: la diferencia entre retardos sucesivos, que la RFC 3393 contempla y de la que iperf3 reporta una media suavizada del tipo definido en la RFC 3550 \cite{rfc3550}, y la variación de dos puntos respecto del retardo mínimo, que es la que emplea la Recomendación ITU-T Y.1541 para fijar sus objetivos. La Figura~\ref{fig:parametros_y1540} ilustra esta diferencia, que el protocolo de medición deberá resolver de manera explícita para que el contraste con los objetivos sea válido. Conviene, por último, distinguir la capacidad de la ruta de la tasa que obtiene una aplicación, pues los protocolos de capa superior con control de flujo en lazo cerrado, como TCP, influyen en la tasa que se experimenta en la capa IP \cite{ituy1540}. A la tasa de datos útiles que recibe la aplicación, descontadas las cabeceras y las retransmisiones, se la denomina \textit{goodput}.

\begin{figure}[htbp]
  \centering
  % Parámetros de Y.1540: retardo en un sentido (IPTD), variación del retardo y pérdida
\begin{tikzpicture}[x=1cm, y=1cm,
  paquete/.style={-{Stealth[length=2mm]}, thick, FigAcento},
  rot/.style={font=\scriptsize}]
  % Líneas de tiempo
  \draw[thick] (0,2.4) node[left, font=\footnotesize] {Origen} -- (11.4,2.4);
  \draw[thick, -{Stealth[length=2mm]}] (0,0) node[left, font=\footnotesize] {Destino} -- (11.6,0)
    node[below left, rot] {tiempo};
  % Paquetes entregados: envío s y retardo x
  \foreach \s/\d/\i in {0.6/1.0/1, 3.0/1.7/2, 5.4/1.2/3}{
    \draw[paquete] (\s,2.4) -- (\s+\d,0);
    \node[rot, above] at (\s,2.4) {P\i};
    \node[rot, fill=white, inner sep=1pt] at (\s+\d/2+0.05,1.2) {$x_{\i}$};
  }
  % Paquete perdido
  \draw[thick, FigAcento, dashed] (7.8,2.4) -- (8.5,1.1);
  \node[rot, above] at (7.8,2.4) {P4};
  \node[font=\footnotesize\bfseries, text=red!70!black] at (8.55,1.0) {$\times$};
  \node[rot, anchor=west] at (8.8,1.0) {Paquete perdido};
  % Separación entre paquetes en el origen y en el destino
  \draw[decorate, decoration={brace, amplitude=4pt}] (0.6,2.95) -- (3.0,2.95)
    node[midway, above=4pt, rot] {$\Delta$};
  \draw[decorate, decoration={brace, amplitude=4pt, mirror}] (1.6,-0.15) -- (4.7,-0.15)
    node[midway, below=4pt, rot] {$\Delta + (x_2 - x_1)$};
  % Definiciones
  \node[rot, align=left, anchor=north west, text width=12.6cm] at (-1.4,-1.0) {%
    \textbf{IPTD:} retardo en un sentido $x_k$ de cada paquete entre los puntos de medición de origen y de destino.\\
    \textbf{Variación entre paquetes sucesivos} (RFC 3393): $x_2 - x_1$; iperf3 reporta su media suavizada (RFC 3550).\\
    \textbf{IPDV de dos puntos} (Y.1540): $v_k = x_k - x_{\min}$; la Recomendación Y.1541 acota su cuantil $1-10^{-3}$.\\
    \textbf{IPLR:} proporción de paquetes transmitidos que no llegan al destino.};
\end{tikzpicture}

  \caption{Retardo en un sentido, variación del retardo y pérdida de paquetes entre dos puntos de medición.}
  \label{fig:parametros_y1540}
  \raggedright
  \elaboradopor{Josué Gabriel Godoy Apolo}
  \fuente{Elaboración propia con base en la Recomendación ITU-T Y.1540 y en las RFC 3393 y 3550 \cite{ituy1540,rfc3393,rfc3550}}
\end{figure}

\subsection{Objetivos de desempeño de la Recomendación ITU-T Y.1541}

Medir, no obstante, no es todavía evaluar: para juzgar un valor es preciso compararlo con un umbral. Este papel lo cumple la Recomendación ITU-T Y.1541, que establece objetivos de desempeño de red para los servicios basados en IP, organizados en clases de calidad de servicio según el tipo de aplicación \cite{ituy1541}. Como muestra la Tabla~\ref{tab:clases_y1541}, las clases 0 y 1, destinadas a aplicaciones en tiempo real como la voz y la videoconferencia, limitan la media del retardo de transferencia a 100~ms y 400~ms, respectivamente, y la variación del retardo a 50~ms; las clases 2 a 4 dejan sin especificar la variación del retardo, y la clase 5 no fija objetivo alguno. Todas las clases con objetivos limitan la tasa de pérdida a $1\times10^{-3}$.

\begin{table}[htbp]
  \centering
  \caption{Clases de calidad de servicio y objetivos de desempeño de la Recomendación ITU-T Y.1541.}
  \label{tab:clases_y1541}
  \footnotesize
  \renewcommand{\arraystretch}{1.2}
  \setlength{\tabcolsep}{4pt}
  \begin{tabularx}{\textwidth}{|c|c|c|c|c|>{\raggedright\arraybackslash}X|}
    \hline
    \textbf{Clase} & \textbf{IPTD} & \textbf{IPDV} & \textbf{IPLR} & \textbf{IPER} & \textbf{Aplicaciones (ejemplos)} \\
    \hline
    0 & 100 ms & 50 ms & $1\times10^{-3}$ & $1\times10^{-4}$ & Tiempo real, sensibles al \textit{jitter} y de alta interacción (voz sobre IP, videoconferencia) \\
    \hline
    1 & 400 ms & 50 ms & $1\times10^{-3}$ & $1\times10^{-4}$ & Tiempo real, sensibles al \textit{jitter} e interactivas (voz sobre IP, videoconferencia) \\
    \hline
    2 & 100 ms & U & $1\times10^{-3}$ & $1\times10^{-4}$ & Datos transaccionales altamente interactivos (señalización) \\
    \hline
    3 & 400 ms & U & $1\times10^{-3}$ & $1\times10^{-4}$ & Datos transaccionales interactivos \\
    \hline
    4 & 1 s & U & $1\times10^{-3}$ & $1\times10^{-4}$ & Solo baja pérdida (transacciones cortas, datos masivos, video en flujo continuo) \\
    \hline
    5 & U & U & U & U & Aplicaciones tradicionales de las redes IP por defecto \\
    \hline
  \end{tabularx}
  \vspace{2pt}
  \raggedright
  \elaboradopor{Josué Gabriel Godoy Apolo}
  \fuente{Traducido y adaptado de la Recomendación ITU-T Y.1541, tablas 1 y 2 \cite{ituy1541}}
  \noindent\textit{Nota:} IPTD, cota superior de la media del retardo de transferencia; IPDV, cota superior del cuantil $1-10^{-3}$ del retardo menos el retardo mínimo; IPLR, cota superior de la probabilidad de pérdida; IPER, cota superior de la tasa de paquetes con error; U, no especificado.
\end{table}

Dos precisiones condicionan el uso de estos objetivos en el presente trabajo. La primera es que se aplican extremo a extremo, entre los puntos de medición que delimitan la red IP, con un intervalo de evaluación sugerido de un minuto. La segunda es que la cota de variación del retardo se refiere al cuantil $1-10^{-3}$, cuya estimación exige al menos mil paquetes por muestra \cite{ituy1541}. Ambas imponen condiciones al diseño de las sesiones de medición, que se fijarán en la metodología.

\subsection{El protocolo TCP y los factores que limitan su desempeño}

La mayor parte del tráfico que percibe el usuario viaja sobre el protocolo de control de transmisión (TCP), cuya especificación vigente recoge la RFC 9293 \cite{rfc9293}. TCP establece cada conexión mediante un intercambio de tres segmentos ---SYN, SYN-ACK y ACK--- que consume un tiempo de ida y vuelta antes de que fluya el primer dato. Durante la transferencia, dos mecanismos regulan la cantidad de datos en tránsito. El control de flujo limita al emisor a la ventana anunciada por el receptor, cuyo campo de 16 bits solo puede superar los 65\,535 bytes mediante la opción de escalado de ventana \cite{rfc7323}. El control de congestión, por su parte, ajusta una ventana de congestión que crece de forma exponencial en el arranque lento y de forma más gradual en la fase de evitación de la congestión, y que se reduce ante la pérdida; la retransmisión rápida se dispara, en particular, al recibir tres acuses de recibo duplicados \cite{rfc5681}. El algoritmo CUBIC, adoptado como predeterminado en las pilas de Linux, Windows y Apple, rige hoy esta dinámica en la mayoría de los equipos cliente \cite{rfc9438}.

De estos mecanismos se desprenden los límites que el análisis deberá identificar. Una conexión no puede transmitir más de una ventana por cada tiempo de ida y vuelta, de modo que su tasa queda acotada por el cociente entre la ventana efectiva y el RTT; para aprovechar la capacidad del cuello de botella de la ruta, la ventana debe igualar al menos el producto ancho de banda-retardo (BDP) \cite{rfc6349}, como muestra la Figura~\ref{fig:ventana_bdp}. Cuando hay pérdida, la tasa sostenida decrece, además, en proporción inversa a la raíz cuadrada de la tasa de pérdida \cite{mathis1997macroscopic}. El medio inalámbrico añade una dificultad propia: TCP interpreta toda pérdida como señal de congestión, de modo que las pérdidas debidas a errores de transmisión en el enlace inalámbrico reducen la ventana sin que la red esté congestionada \cite{balakrishnan1997comparison}. Del mismo modo, el reordenamiento de paquetes genera acuses duplicados que pueden desencadenar retransmisiones innecesarias. Retransmisiones, reordenamiento y limitación por ventana son, por tanto, las huellas que la disección de las sesiones debe buscar para atribuir correctamente una tasa baja.

\begin{figure}[htbp]
  \centering
  % Tasa de una conexión TCP según la ventana efectiva: límite de la ventana y límite de la ruta
\begin{tikzpicture}[x=1cm, y=1cm, rot/.style={font=\scriptsize}]
  % Ejes
  \draw[-{Stealth[length=2mm]}, thick] (0,0) -- (10.4,0);
  \node[rot, anchor=north east] at (10.4,-0.45) {Ventana efectiva $W$};
  \draw[-{Stealth[length=2mm]}, thick] (0,0) -- (0,4.6) node[above, rot] {Tasa TCP $T$};
  % Curva: crece con la ventana hasta el producto ancho de banda-retardo
  \draw[FigAcento, very thick] (0,0) -- (5,3.5) -- (9.8,3.5);
  \draw[densely dashed, gray!70!black] (5,0) node[below, rot] {$W = \mathrm{BDP} = \mathrm{BB} \times \mathrm{RTT}$} -- (5,3.5);
  \draw[densely dashed, gray!70!black] (0,3.5) node[left, rot] {BB} -- (5,3.5);
  % Regiones
  \node[rot, align=center] at (1.7,2.9) {Limitada por la ventana\\$T = W/\mathrm{RTT}$};
  \node[rot, align=center] at (7.4,4.05) {Limitada por la ruta\\$T = \mathrm{BB}$};
  % Leyenda
  \node[rot, anchor=north west, text width=10.4cm, align=left] at (0,-1.0)
    {BB: ancho de banda del cuello de botella de la ruta; RTT: tiempo de ida y vuelta; BDP: producto ancho de banda-retardo.};
\end{tikzpicture}

  \caption{Tasa de una conexión TCP según la ventana efectiva y el producto ancho de banda-retardo.}
  \label{fig:ventana_bdp}
  \raggedright
  \elaboradopor{Josué Gabriel Godoy Apolo}
  \fuente{Elaboración propia con base en la RFC 6349 \cite{rfc6349}}
\end{figure}

\subsection{Metodología de la RFC 6349}

Que la tasa TCP dependa de tantos factores explica que su medición requiera un procedimiento normalizado. La RFC 6349 establece un marco práctico para medir el rendimiento TCP extremo a extremo en redes IP gestionadas, en tres pasos sucesivos: identificar la unidad máxima de transmisión de la ruta, establecer el tiempo de ida y vuelta de referencia y el ancho de banda del cuello de botella, y ejecutar pruebas de rendimiento TCP con una o varias conexiones \cite{rfc6349}. La medición se realiza en el estado de equilibrio de la conexión, una vez superado el arranque lento, y sus resultados se expresan mediante tres métricas: la relación entre el tiempo de transferencia real y el ideal (\textit{TCP Transfer Time Ratio}); la eficiencia TCP, que expresa la proporción de los bytes transmitidos que no debieron retransmitirse; y el porcentaje de retardo por búfer, que cuantifica el aumento del RTT durante la transferencia respecto del valor de referencia. La propia RFC aclara que su metodología no pretende ofrecer un diagnóstico detallado de los problemas del extremo o de la red, de ahí que el presente trabajo la complemente con la disección pasiva. Concebida, además, para redes gestionadas, en las que el operador controla ambos extremos, su aplicación desde el equipo de un usuario hacia destinos de referencia constituye una adaptación que el protocolo de medición deberá declarar.

Distinto es el caso de la RFC 2544, que define una metodología de evaluación comparativa de dispositivos de interconexión en laboratorio \cite{rfc2544}. Sus pruebas sobrecargan deliberadamente los recursos del dispositivo bajo prueba para determinar su capacidad, razón por la cual la RFC 6815 restringe su aplicación a entornos de prueba aislados y advierte que su uso en redes en producción perjudicaría el tráfico de los usuarios \cite{rfc6815}. En este trabajo, la RFC 2544 se emplea, por consiguiente, solo como referencia conceptual de terminología y de procedimiento.

\subsection{Unidad máxima de transmisión y su descubrimiento}

La unidad máxima de transmisión (MTU) es el tamaño del mayor paquete que un enlace puede transportar sin fragmentarlo; en Ethernet es de 1500 bytes, pero los túneles y algunos enlaces del proveedor la reducen. La MTU de una ruta es, por tanto, la menor de las MTU de sus enlaces. Para descubrirla, el mecanismo clásico (PMTUD) envía paquetes con la indicación de no fragmentar y espera que el enrutador incapaz de reenviarlos responda con un mensaje ICMP que informe la MTU admisible, tanto en IPv4 \cite{rfc1191} como en IPv6 \cite{rfc8201}. El mecanismo es frágil precisamente por depender de esos mensajes: cuando un cortafuegos los descarta, los paquetes grandes se pierden en silencio y las conexiones se detienen sin error aparente, un fenómeno que Luckie y Stasiewicz midieron en una proporción no despreciable de servidores \cite{luckie2010pmtud}. Como alternativa, el descubrimiento en la capa de empaquetamiento (PLPMTUD) sondea la ruta con paquetes de tamaño creciente y no depende de los mensajes ICMP \cite{rfc4821}. Verificar la MTU efectiva de la ruta y la operatividad de su descubrimiento constituye, además, el primer paso de la metodología de la RFC 6349.

\subsection{Fases del establecimiento de un servicio}

El retardo que percibe un usuario al abrir un servicio no es un bloque indivisible, sino la suma de fases que se suceden en un orden fijo (Figura~\ref{fig:fases_servicio}). Al asociarse a la red, el equipo obtiene su configuración IP mediante el intercambio DHCP de descubrimiento, oferta, solicitud y confirmación \cite{rfc2131}. Para cada servicio, el primer paso es la \textbf{resolución del nombre}: el equipo consulta a un servidor DNS recursivo, que responde desde su caché o tras consultar a los servidores autoritativos \cite{rfc1034}. Le sigue el \textbf{establecimiento de la conexión TCP}, que consume un tiempo de ida y vuelta \cite{rfc9293}, y, en los servicios cifrados, la \textbf{negociación TLS}, que en la versión 1.3 completa el intercambio de claves y la autenticación del servidor en un tiempo de ida y vuelta adicional \cite{rfc8446}. Solo entonces el cliente envía su solicitud, espera el primer byte de la respuesta ---tiempo en el que se suma el procesamiento del servidor--- y recibe la \textbf{transferencia} del contenido.

Cada fase tiene causas de degradación propias: una resolución lenta apunta al servicio de nombres; un establecimiento lento, al retardo de la ruta; una espera prolongada del primer byte, al servidor; y una transferencia lenta, a la capacidad disponible o al control de congestión. Descomponer el retardo percibido en estas fases es, por consiguiente, la condición para atribuir correctamente la degradación \cite{sundaresan2013web}.

\begin{figure}[htbp]
  \centering
  % Fases del establecimiento de un servicio web cifrado: DNS, TCP, TLS 1.3, primer byte y transferencia
\begin{tikzpicture}[x=1cm, y=1.3cm,
  msj/.style={-{Stealth[length=1.8mm]}, thick, FigAcento},
  etq/.style={font=\scriptsize, fill=white, inner sep=1pt},
  rot/.style={font=\scriptsize}]
  % Participantes y líneas de vida
  \foreach \x/\n/\t in {0/cli/Cliente, 4.5/dns/Resolutor DNS, 9.0/srv/Servidor}{
    \node[caja, minimum width=2.0cm] (\n) at (\x,0.5) {\t};
    \draw[gray!60, thick] (\x,0.2) -- (\x,-5.8);
  }
  % Resolución DNS
  \draw[msj] (0,-0.3) -- node[etq] {Consulta} (4.5,-0.7);
  \draw[msj] (4.5,-0.9) -- node[etq] {Respuesta} (0,-1.3);
  % Conexión TCP
  \draw[msj] (0,-1.6) -- node[etq, pos=0.25] {SYN} (9.0,-2.0);
  \draw[msj] (9.0,-2.0) -- node[etq, pos=0.25] {SYN-ACK} (0,-2.4);
  % Negociación TLS 1.3
  \draw[msj] (0,-2.4) -- node[etq, pos=0.25] {ACK + ClientHello} (9.0,-2.8);
  \draw[msj] (9.0,-2.8) -- node[etq, pos=0.25] {ServerHello \dots{} Finished} (0,-3.2);
  % Solicitud y primer byte
  \draw[msj] (0,-3.2) -- node[etq, pos=0.25] {Finished + solicitud HTTP} (9.0,-3.6);
  \node[rot, anchor=east, text=gray!70!black] at (8.85,-3.8) {procesamiento};
  \draw[msj] (9.0,-4.0) -- node[etq, pos=0.35] {Primer byte de la respuesta} (0,-4.4);
  % Transferencia
  \draw[msj] (9.0,-4.6) -- (0,-5.0);
  \draw[msj] (9.0,-4.9) -- node[etq, pos=0.25] {Datos} (0,-5.3);
  \draw[msj] (9.0,-5.2) -- (0,-5.6);
  % Fases medidas en el cliente
  \foreach \a/\b/\t in {-1.3/-0.3/Resolución DNS, -2.4/-1.6/Conexión TCP, -3.2/-2.4/Negociación TLS,
      -4.4/-3.2/Espera del primer byte, -5.6/-4.4/Transferencia}{
    \draw[decorate, decoration={brace, amplitude=4pt}] (-0.25,\a) -- (-0.25,\b)
      node[midway, left=5pt, rot] {\t};
  }
  % Correspondencia con los tiempos de curl
  \node[rot, anchor=north west, text width=12.4cm, align=left] at (-3.1,-6.0)
    {curl registra el instante acumulado en que termina cada fase: \texttt{time\_namelookup}, \texttt{time\_connect}, \texttt{time\_appconnect}, \texttt{time\_starttransfer} y \texttt{time\_total}.};
\end{tikzpicture}

  \caption{Fases del establecimiento de un servicio web cifrado con TLS 1.3.}
  \label{fig:fases_servicio}
  \raggedright
  \elaboradopor{Josué Gabriel Godoy Apolo}
  \fuente{Elaboración propia con base en las RFC 1034, 9293 y 8446 \cite{rfc1034,rfc9293,rfc8446}}
\end{figure}

\subsection{Acceso inalámbrico IEEE 802.11}

En la mayoría de los casos, el primer tramo de la ruta es un enlace inalámbrico IEEE 802.11 \cite{ieee80211}, cuyo comportamiento difiere del cableado en un aspecto esencial: el medio es compartido y su calidad varía con el lugar y con el momento. La \textbf{intensidad de la señal recibida} (RSSI), expresada en dBm, mide la potencia con que llega la señal del punto de acceso; la \textbf{relación señal-ruido} (SNR), en dB, la compara con el ruido de fondo del canal. De la SNR depende el \textbf{esquema de modulación y codificación} (MCS) que el enlace puede sostener: a mayor SNR, modulaciones con más bits por símbolo y codificaciones menos redundantes. La combinación del MCS, el ancho de canal y el número de flujos espaciales determina la \textbf{tasa nominal} de la capa física, que fijan las sucesivas enmiendas del estándar, como 802.11n (Wi-Fi 4), 802.11ac (Wi-Fi 5) y 802.11ax (Wi-Fi 6); el \textit{goodput} que obtiene el usuario es siempre una fracción de esa tasa, por la sobrecarga del protocolo y la contienda por el medio.

Cada trama de datos dirigida a un único destino debe, además, ser confirmada por el receptor, y la que no se confirma se retransmite. La \textbf{tasa de reintentos} revela así las colisiones y la interferencia que la intensidad de la señal, por sí sola, no muestra. Esa interferencia es especialmente probable en la banda de 2,4~GHz, cuyos canales, separados apenas 5~MHz y de 20 a 22~MHz de ancho, se superponen entre sí, de modo que solo tres de ellos ---el 1, el 6 y el 11--- pueden operar en un mismo espacio sin interferirse (Figura~\ref{fig:canales_24}). La mayoría de estos indicadores es expuesta por el controlador de la interfaz inalámbrica del propio cliente, aunque su disponibilidad varía según el controlador; ello permite registrarlos junto con cada medición sin acceso a la infraestructura.

\begin{figure}[htbp]
  \centering
  % Canales 1 a 11 de la banda de 2,4 GHz (fc = 2407 + 5n MHz, ancho aprox. 22 MHz)
% Eje x: MHz por encima de 2400
\begin{tikzpicture}[x=0.15cm, y=1cm]
  % Canales superpuestos
  \foreach \n in {2,3,4,5,7,8,9,10}{
    \pgfmathsetmacro{\fc}{7+5*\n}
    \draw[gray!65, dashed] plot[domain=-11:11, samples=41] ({\fc+\x},{1.8*cos(\x*180/22)});
    \node[font=\scriptsize, text=gray!70!black, above] at (\fc,1.8) {\n};
  }
  % Canales sin solapamiento
  \foreach \n in {1,6,11}{
    \pgfmathsetmacro{\fc}{7+5*\n}
    \fill[FigAcento, opacity=0.18] plot[domain=-11:11, samples=41] ({\fc+\x},{1.8*cos(\x*180/22)}) -- cycle;
    \draw[FigAcento, very thick] plot[domain=-11:11, samples=41] ({\fc+\x},{1.8*cos(\x*180/22)});
    \node[font=\scriptsize\bfseries, text=FigAcento, above] at (\fc,1.8) {\n};
  }
  % Eje de frecuencia
  \draw[-{Stealth[length=2mm]}] (-2,0) -- (90,0);
  \foreach \f/\etiqueta in {0/2400, 12/2412, 37/2437, 62/2462, 83.5/{2483,5}}{
    \draw (\f,0) -- (\f,-0.12) node[below, font=\scriptsize] {\etiqueta};
  }
  \node[font=\scriptsize, anchor=north east] at (90,-0.5) {Frecuencia (MHz)};
  % Leyenda
  \fill[FigAcento, opacity=0.18] (2,-1.3) rectangle (6,-1.0);
  \draw[FigAcento, very thick] (2,-1.3) rectangle (6,-1.0);
  \node[font=\scriptsize, anchor=west] at (6.5,-1.15) {Canales sin solapamiento (1, 6 y 11)};
  \draw[gray!65, dashed] (42,-1.15) -- (48,-1.15);
  \node[font=\scriptsize, anchor=west] at (48.5,-1.15) {Canales que se superponen};
\end{tikzpicture}

  \caption{Canales 1 a 11 de la banda de 2,4~GHz y canales sin solapamiento.}
  \label{fig:canales_24}
  \raggedright
  \elaboradopor{Josué Gabriel Godoy Apolo}
  \fuente{Elaboración propia con base en IEEE 802.11 \cite{ieee80211}}
\end{figure}

\subsection{Línea base y tratamiento estadístico}

Las mediciones, sin embargo, solo adquieren significado cuando se las compara. Una línea base de desempeño es el registro del comportamiento normal de la red en condiciones conocidas, que sirve de referencia para juzgar cualquier medición posterior \cite{oppenheimer2011topdown}. Su valor no reside en un número aislado, sino en la reproducibilidad del procedimiento que lo produjo: una línea base debe especificar qué se midió, dónde, cuándo, con qué herramienta y bajo qué carga. En el presente trabajo, esa especificación adopta la forma de un diseño muestral por punto de asociación y franja horaria.

Como el desempeño de una red de campus presenta distribuciones asimétricas y de cola larga \cite{sui2016wifiseer}, la línea base se describe con medidas robustas ---la mediana y los percentiles para la tendencia central y las colas, el rango intercuartílico para la dispersión---, que la media y la desviación estándar resumirían mal. Por la misma razón, la relación entre los parámetros de capas superiores y los indicadores del acceso inalámbrico se estima mediante el coeficiente de correlación por rangos de Spearman, que no supone una relación lineal ni una distribución normal de las variables \cite{hollander2014nonparametric}.

\subsection{Ciclo de vida de red PPDIOO}

Si los apartados anteriores describen qué se mide y cómo, queda por situar el trabajo en el tiempo de la red. El ciclo PPDIOO divide la vida de una red en seis fases: preparación, planificación, diseño, implementación, operación y optimización \cite{oppenheimer2011topdown}. Las dos primeras se ocupan de comprender los requisitos de la organización y de caracterizar la red existente antes de cualquier decisión de diseño, y es en ellas donde se desarrolla el presente trabajo (Figura~\ref{fig:ppdioo}). La preparación delimita los puntos de asociación, la metodología de medición y los parámetros a registrar; la planificación elabora el protocolo, ejecuta la campaña y construye la línea base. Sus resultados alimentan las fases posteriores, cuya ejecución corresponderá a la institución.

\begin{figure}[htbp]
  \centering
  % Ciclo de vida PPDIOO; se resaltan las fases que aborda el trabajo
\begin{tikzpicture}[
  fase/.style={caja, fill=white, font=\footnotesize, minimum width=2.9cm, minimum height=0.9cm},
  abordada/.style={fill=FigSuave, font=\footnotesize\bfseries}]
  \foreach \ang/\nombre/\estilo [count=\i] in {
      90/Preparación/abordada, 30/Planificación/abordada, -30/Diseño/,
      -90/Implementación/, -150/Operación/, 150/Optimización/}{
    \node[fase, \estilo] (f\i) at ({3.9*cos(\ang)},{2.3*sin(\ang)}) {\nombre};
  }
  \foreach \i/\j in {1/2, 2/3, 3/4, 4/5, 5/6, 6/1}{
    \draw[flecha] (f\i) to[bend left=18] (f\j);
  }
  \node[align=center, font=\footnotesize\itshape, text=FigAcento] at (0,0) {Ciclo de vida\\de la red};
  % Leyenda
  \node[caja, minimum width=0.6cm, minimum height=0.35cm] (ley) at (-2.6,-3.4) {};
  \node[anchor=west, font=\scriptsize] at (ley.east) {Fases en las que se desarrolla el presente trabajo};
\end{tikzpicture}

  \caption{Ciclo de vida de red PPDIOO y fases en las que se desarrolla el presente trabajo.}
  \label{fig:ppdioo}
  \raggedright
  \elaboradopor{Josué Gabriel Godoy Apolo}
  \fuente{Elaboración propia con base en Oppenheimer \cite{oppenheimer2011topdown}}
\end{figure}

\subsection{Clasificación de hallazgos y análisis de riesgo}

Caracterizar, finalmente, no consiste en enumerar anomalías, sino en jerarquizarlas. Una degradación que afecta a la franja de mayor ocupación de un espacio de uso masivo no tiene el mismo peso que otra que aparece de forma esporádica en un punto secundario. Para establecer esa jerarquía, el trabajo adopta el enfoque de evaluación de riesgos propuesto por el NIST en su publicación especial 800-30, según el cual el riesgo se determina a partir de la probabilidad de que una condición desfavorable se materialice y de la magnitud de su impacto, valoradas mediante escalas cualitativas o semicuantitativas \cite{nist80030}. La Figura~\ref{fig:matriz_riesgo} reproduce la escala con la que dicha publicación combina ambas dimensiones en cinco niveles de riesgo. Trasladada al presente trabajo, la probabilidad se aproxima por la frecuencia con que la condición se observa en el diseño muestral, y el impacto, por su efecto sobre la experiencia del usuario. La matriz de hallazgos resultante clasifica cada condición por criticidad y constituye el puente entre la caracterización y las recomendaciones técnicas, pues determina el orden en que convendría atenderlas. Con ella se cierra el marco conceptual sobre el que se apoyará la metodología expuesta en el capítulo siguiente.

\begin{figure}[htbp]
  \centering
  % Nivel de riesgo según probabilidad e impacto (NIST SP 800-30 Rev. 1, tabla I-2)
% Filas de abajo hacia arriba: probabilidad muy baja ... muy alta
% Valores: 1 muy bajo, 2 bajo, 3 moderado, 4 alto, 5 muy alto
\begin{tikzpicture}[x=2.0cm, y=0.8cm]
  \definecolor{riesgo1}{HTML}{C6E5B3}
  \definecolor{riesgo2}{HTML}{E4EFA6}
  \definecolor{riesgo3}{HTML}{FFE699}
  \definecolor{riesgo4}{HTML}{F8B98A}
  \definecolor{riesgo5}{HTML}{F08A80}
  \foreach \fila [count=\r] in {{1,1,1,2,2},{1,2,2,2,3},{1,2,3,3,4},{1,2,3,4,5},{1,2,3,4,5}}{
    \foreach \nivel [count=\c] in \fila {
      \fill[riesgo\nivel] (\c-1,\r-1) rectangle (\c,\r);
      \draw[white, line width=1.2pt] (\c-1,\r-1) rectangle (\c,\r);
      \node[font=\scriptsize] at (\c-0.5,\r-0.5)
        {\ifcase\nivel\or Muy bajo\or Bajo\or Moderado\or Alto\or Muy alto\fi};
    }
  }
  % Rótulos de probabilidad (filas) e impacto (columnas)
  \foreach \t [count=\r] in {Muy baja, Baja, Moderada, Alta, Muy alta}{
    \node[font=\scriptsize, anchor=east] at (-0.05,\r-0.5) {\t};
  }
  \foreach \t [count=\c] in {Muy bajo, Bajo, Moderado, Alto, Muy alto}{
    \node[font=\scriptsize, anchor=north] at (\c-0.5,-0.05) {\t};
  }
  \node[font=\footnotesize\bfseries, text=FigAcento, rotate=90] at (-0.95,2.5) {Probabilidad};
  \node[font=\footnotesize\bfseries, text=FigAcento] at (2.5,-1.05) {Impacto};
\end{tikzpicture}

  \caption{Nivel de riesgo según la combinación de probabilidad e impacto.}
  \label{fig:matriz_riesgo}
  \raggedright
  \elaboradopor{Josué Gabriel Godoy Apolo}
  \fuente{Adaptado de NIST SP 800-30 Rev.~1, tabla I-2 \cite{nist80030}}
\end{figure}
